% D04-001. Согласованная композиция 05.10, продолжение 09.10.2026.
\chapter{Методы анализа осциллограмм и распознавания событий}
\label{ch:statistical_analysis}
\label{ch:oscillogram_analysis}

\section{Постановка исследования}
\label{sec:ch04_intro}

Подготовка реальных осциллограмм создаёт основу для исследования
методов анализа аварийных и переходных процессов. Статистическое
использование данных рассмотрено в главе~\ref{ch:dataset_preparation}.
В настоящей главе рассматриваются представления сигналов
и нейросетевые способы их обработки. Особое внимание уделяется
включению электротехнических соотношений в архитектуру модели
и различиям между распознаванием события и решением
измерительного органа.

Для БАВР необходима информация о развитии процесса и состоянии
сети, по которой определяются условия восстановления питания.
Распознавание ОЗЗ, короткого замыкания, коммутаций и других явлений позволяет
описать последовательность событий в зарегистрированном материале.
Решение о разрешении либо блокировке переключения дополнительно
требует оценки множества условий работы.
Подобные подходы могут применяться также и для анализа
и предметной проверкой алгоритмов РЗА.
Поэтому исследование моделей обработки осциллограмм является важной задачей.

Далее излагаются принципы физически мотивированных архитектур.
Их развитие связано с выбором признаков, использованием
комплексных представлений и исследованием детекции событий.
Эти направления рассматриваются в единой связи с подготовкой
данных и задачами совершенствования БАВР.

% D04-001. AI-01 P0064--P0092, уточнения по локальным реализациям.
\section{Физически мотивированные нейросетевые архитектуры}
\label{sec:ch4_physkan_architectures}

При анализе осциллограмм нейронная сеть должна выделять признаки
процесса из сигналов токов и напряжений. Выбор архитектуры определяет,
какие связи между этими сигналами задаются заранее, а какие должны
быть установлены в ходе обучения. Для задач РЗА представляет интерес
введение в вычисления известных электротехнических соотношений.
Они дают исходную структуру для обработки сигналов, а обучаемая часть
позволяет учитывать особенности зарегистрированных процессов.

В настоящей работе обозначение ФизКАН относится к семейству
модификаций нейронных сетей, объединяющих физические операции
и слои KAN. Развитие этого подхода
рассматривается на примере вещественной модели PhysicsKAN,
комплексной cPhysicsKAN и модели rPhysicsKAN с управлением
амплитудными признаками по фазовым соотношениям.\footnote{Евдаков А.Е.
Архитектурная модернизация интеллектуальных алгоритмов релейной защиты.
Внедрение законов электротехники в граф вычислений.}
Эти варианты отражают различные способы включения физической
информации в обработку осциллограммы.

\subsection{Базовые модели обработки сигналов}

Для сопоставления способов обработки рассматриваются полносвязная,
свёрточная и остаточная архитектуры. В модели SimpleMLP входное
окно преобразуется в вектор признаков, который поступает
в последовательность полносвязных слоёв. 
(Примечание автора: Есть мысли о том, чтобы картинки тут привести поясняющие схему, 
если у меня в оригинальных работах есть. Убрать потом это, в случае чего будет не сложно, 
а пока что хочу максимально подробно переписывать статьи. 
Кстати, об этом опять же можешь сделать пометку, она будет полезна для будущего текста
Причём... Можно подумать о том, чтобы нарисовать какую-то другую картинку поясняющую, как я делал в будущих работах, 
я думаю современные ИИ модели могут помочь сдлеать её тут, надо будет только инструкцию хорошую написать для картинки,
если тебе сложно. И да, лучше через инструкцию сделать это, а не тебе самой выдумывать, причём я хочу тут иснтрукцию
поясняющую проблему для энергетики, где у меня подаются и ТОКИ и НАПРЯЖЕНИЯ и по времени распределены и они путаются.
Да, давай попробуем через картинки такие новые.
Хмм... В статье есть картинки, но они довольно сложны, и не факт, что отражают физическую проблему).
(Примечание автора: И да, ещё раз отмечу, давай побольше текста пока что из оригинальных работ писать сюда.
Да, не весь и не всё, но важное, люди из энергетики многое базовое могу не понимать.).
Взаимосвязи компонентов определяются обучаемыми весами. Порядок отсчётов сохраняется
в положении элементов вектора, однако локальная временная структура
не задаётся специальным механизмом архитектуры. Такая модель
служит исходным вариантом для сравнения более структурированных
способов обработки.

В модели SimpleCNN одномерные свёртки обрабатывают соседние отсчёты
и используют общие коэффициенты в разных положениях окна.
Это позволяет выделять локальные изменения формы сигнала
одним набором обучаемых фильтров. Последовательность слоёв
объединяет локальные признаки в описание рассматриваемого процесса.
Физический смысл каждого фильтра заранее не фиксируется,
поэтому его связь с конкретным измерительным признаком
требует отдельного анализа.

В модели ResNet1D применяются остаточные связи. Выход блока
суммируется с его входом после необходимого согласования
размерностей. Это сохраняет путь передачи исходной информации
и позволяет обучать последовательность преобразований.
Увеличение глубины и числа каналов меняет объём вычислений
и количество параметров. Поэтому архитектура оценивается
совместно с её конкретной конфигурацией и представлением входа.
Само название модели не определяет ни качество распознавания,
ни возможность её применения в микропроцессорном устройстве.

(Примечание автора: Ну да, как-то очень кратко... Мне-то понятно, но я всё это и написал ранее.)

\subsection{Сети Колмогорова--Арнольда}

В сетях Колмогорова--Арнольда обучаемые одномерные функции
располагаются на связях между компонентами соседних слоёв.
В отличие от полносвязного слоя с фиксированной функцией активации,
обучение изменяет форму этих зависимостей. В рассматриваемых моделях
используются функции на основе сплайнов.\footnote{Liu Z. et al.
KAN: Kolmogorov--Arnold Networks. 2025.
\url{https://arxiv.org/abs/2404.19756v5}.}
Для входного вектора $\boldsymbol{x}$ преобразование слоя
можно записать в виде
\begin{equation}
    y_j = \sum_{i=1}^{n} f_{ji}(x_i),
    \label{eq:ch4_kan_layer}
\end{equation}
где $f_{ji}$ обозначает обучаемую функцию на соответствующей связи.
Состав и ширина последовательности слоёв задают сложность модели.

Модель SimpleKAN использует такие преобразования для обработки
входного вектора. Модель ConvKAN сочетает свёрточную обработку
с нелинейными преобразованиями на основе KAN. Временная структура
входа при этом используется для выделения локальных признаков,
а форма обучаемых функций может быть рассмотрена отдельно.

Возможность построить график функции на связи полезна для анализа
модели. Однако наличие такого графика само по себе не устанавливает
электротехнический смысл соответствующего признака. Внутренние
переменные могут представлять обученные комбинации сигналов.
Для содержательного объяснения необходима связь этих переменных
с входными величинами и поведением модели на известных режимах.

(Примечание автора: Из текста не очень понятно, чем это лучше предыдущего?
Я закладывал во многом то, чтобы могли возникать сильно нелинейные связи например от тока и т.д. Ну и да, побольше бы описания…)

\subsection{Вещественные физические операции}

В архитектуре PhysicsKAN введены операции попарного умножения
и деления входных каналов. При согласованном расположении токов
и напряжений эти операции формируют признаки, связанные
с их произведениями и отношениями. Для мгновенных значений
однофазных сигналов произведение $u(t)i(t)$ соответствует
мгновенной мощности. Отношение $i(t)/u(t)$ характеризует
взаимосвязь сигналов, но при переходном процессе не становится
автоматически оценкой постоянной проводимости элемента сети.

При делении требуется учитывать малые значения знаменателя.
Для вещественного сигнала используется ограничение его абсолютной
величины с сохранением знака. Такой приём предотвращает деление
на ноль, но при близком к нулю напряжении отношение всё ещё может
принимать большие значения. Поэтому вычислительная защита
и физическая пригодность признака рассматриваются отдельно.

Сформированные каналы проходят нормализацию и объединяются
с исходными каналами. Последующая обработка сохраняет доступ
и к исходным сигналам, и к результатам физических операций.
Обучаемая часть определяет, какие из них существенны
для рассматриваемой задачи. Арифметические операции включены
в общий граф вычислений, через который проходит обучение.

В рассматриваемом базовом варианте физические операции выполняются
при входной обработке, перед последовательностью KAN-слоёв.
Их применение к признакам на последующих уровнях представляет
другое архитектурное развитие. После обучаемых преобразований
скрытые каналы могут утрачивать непосредственный смысл токов
и напряжений. Поэтому произведения и отношения таких каналов
нуждаются в собственной интерпретации.

\subsection{Комплексные признаки и фазовые соотношения}

При использовании гармонических составляющих существенны
не только амплитуды, но и фазовые соотношения сигналов.
В архитектуре cPhysicsKAN комплексные величины представлены
парами амплитуды и фазы. Порядок каналов согласован с этим
представлением, а операции выполняются одновременно
над обеими компонентами каждой пары.

Для величин $z_1=A_1e^{\mathrm{j}\varphi_1}$
и $z_2=A_2e^{\mathrm{j}\varphi_2}$ произведение и отношение
имеют вид
\begin{equation}
\begin{aligned}
    z_1z_2 &= A_1A_2e^{\mathrm{j}(\varphi_1+\varphi_2)},\\
    \frac{z_1}{z_2} &= \frac{A_1}{A_2}
        e^{\mathrm{j}(\varphi_1-\varphi_2)}, \qquad A_2>0.
\end{aligned}
\label{eq:ch4_complex_operations}
\end{equation}
В вычислениях знаменатель ограничивается положительной малой
величиной. При исчезновении сигнала такая защита обеспечивает
определённость операции, но не восстанавливает физически
достоверную фазу отсутствующего напряжения.

Комплексное произведение следует отличать от комплексной мощности.
При принятом определении векторов тока и напряжения мощность
вычисляется как $\underline{S}=\underline{U}\underline{I}^{*}$,
где звёздочка обозначает комплексное сопряжение.
Поэтому для неё используется разность фаз напряжения и тока.
Обычное произведение в формуле~\eqref{eq:ch4_complex_operations}
содержит сумму фаз. Признаки, полученные таким умножением,
нельзя без дополнительного преобразования отождествлять
с активной и реактивной мощностью.

Отношение $\underline{I}/\underline{U}$ соответствует комплексной
проводимости при согласованном выборе измеряемых величин.
Для физического толкования результатов важны порядок каналов,
нормализация и выбранная гармоника. Если электрические сигналы
заменены их обучаемыми комбинациями, операция сохраняет
математический смысл, а электротехническое название признака
уже требует отдельного обоснования.

Комплексные дополнительные признаки объединяются с исходным
представлением. В рассматриваемой реализации нормализация
результатов физических операций применяется к амплитудам,
а фазы сохраняются отдельными компонентами. Это поддерживает
их различное назначение в последующей обработке.

\subsection{Управление амплитудными признаками по фазе}

Архитектура rPhysicsKAN развивает комплексное представление
за счёт разделения обработки амплитуд и фаз.
После формирования дополнительных признаков амплитудная
и фазовая компоненты поступают в отдельные KAN-ветви.
Ещё одна ветвь получает фазовую информацию и формирует
коэффициенты управления амплитудными признаками.

Пусть $\boldsymbol{a}$ и $\boldsymbol{\varphi}$ обозначают
векторы входных амплитуд и фаз после подготовки признаков.
Обучаемые преобразования амплитудной и управляющей ветвей
обозначим $F_a$ и $F_g$. В рассматриваемой реализации
управляемая амплитудная ветвь имеет вид
\begin{equation}
    \widetilde{\boldsymbol{a}} =
    \operatorname{softplus}\bigl(F_a(\boldsymbol{a})\bigr)
    \odot \sigma\bigl(F_g(\boldsymbol{\varphi})\bigr),
    \label{eq:ch4_relay_gate}
\end{equation}
где $\odot$ обозначает поэлементное умножение,
$\sigma(x)=1/(1+e^{-x})$,
а $\operatorname{softplus}(x)=\ln(1+e^x)$.
Сигмоидальная функция ограничивает коэффициенты управления
диапазоном от нуля до единицы. Преобразование softplus
обеспечивает неотрицательное значение управляемых признаков.
Полученный результат объединяется с выходом фазовой ветви
и поступает в последующую часть классификатора.

Такая структура задаёт зависимость передачи амплитудной
информации от фазовых соотношений. По своему назначению
она близка к использованию угла между электрическими
величинами в направленном измерительном органе.
При этом обучаемый коэффициент управления не равен
заранее заданной функции $\cos\varphi$ и не обеспечивает
сам по себе тождественность классическому РНМ.
Его характеристика определяется обучением и должна
исследоваться на соответствующих режимах.

Для БАВР такая архитектурная идея интересна возможностью
совместно учитывать амплитудные признаки и направление,
связанное с фазовыми соотношениями. Классификация событий
в осциллограмме и решение о разрешении либо блокировке
переключения имеют разные целевые выходы. Поэтому перенос
архитектуры на задачу РНМ предполагает собственную разметку,
условия применимости и проверку поведения органа во времени.

\subsection{Место архитектурного подхода в исследовании}

Развитие моделей последовательно добавляет структуру
в обработку сигналов. Свёртки используют локальные связи
между отсчётами. KAN-слои дают обучаемые нелинейные зависимости
на связях. Физические операции добавляют произведения
и отношения, а фазовое управление позволяет изменять
передачу амплитудной информации в зависимости от угловых
признаков. Каждый из этих механизмов задаёт определённый
способ использования входных данных.

При экспериментальном сопоставлении необходимо одновременно
учитывать представление сигналов, состав меток, размер модели
и способ обработки окна. Число параметров характеризует
объём обучаемой части, но не заменяет измерения времени
вычислений. Время обработки готового окна также отличается
от времени получения необходимых отсчётов и задержки
решения в потоке. Эти различия существенны для оценки
возможности применения модели в быстродействующей автоматике.

Подготовленные в главе~\ref{ch:dataset_preparation} данные
дают основу для исследования этих моделей на реальных
процессах. Объяснение архитектуры определяет, каким образом
используются сигналы, а сравнение на согласованных выборках
позволяет оценить результат этого выбора. Последующая
проверка конкретного измерительного органа связывает
нейросетевой подход с совершенствованием алгоритмов БАВР.


% Исторический план DISS-005 сохранён дословно, без вывода в PDF.
\iffalse
% ==============================================================================
% ГЛАВА 4 (DISS-005-B2054 .. DISS-005-B2063)
% Статистический анализ функционирования БАВР на основе реальных осциллограмм
% (Планируемая)
% ==============================================================================

% DISS-005-B2054
\chapter{Статистический анализ функционирования БАВР на основе реальных осциллограмм (Планируемая)}
\label{ch:statistical_analysis}

% DISS-005-B2055
\section{Введение}
\label{sec:ch04_intro}

% DISS-005-B2056
\textbf{Статус:} \textcolor{red}{В планах.}

% DISS-005-B2057
\textbf{Содержание (предполагаемое):}

% DISS-005-B2058
Использование датасета из Главы 3 для проведения глубокого статистического анализа.

% DISS-005-B2059
Анализ временных характеристик коммутационного оборудования, задействованного в циклах БАВР (реальные времена включения/отключения выключателей).

% DISS-005-B2060
Исследование уровней перенапряжений при ОЗЗ в реальных сетях \textcolor{red}{и их влияния на работу БАВР} (Предварительный набросок уже есть).

% DISS-005-B2061
\dots

% DISS-005-B2062

% DISS-005-B2063

\clearpage

\fi
\clearpage
